\documentclass[10pt,a4paper]{amsart}
\usepackage[margin=19mm]{geometry}
\usepackage{amsmath,amssymb,mathtools}
\usepackage[scr=rsfs,cal=euler]{mathalfa}
\usepackage{xfrac}
\usepackage[unicode,hidelinks,bookmarks=false]{hyperref}
\newcommand{\ie}{i.e.\ }
\newcommand{\cf}{cf.\ }
\newcommand{\resp}{respectively\ }
\newcommand{\Subsection}{\subsection}
\providecommand{\nobreakdash}{\nobreak\hbox{-}}
\setlength{\emergencystretch}{2em}
\multlinegap0pt
\usepackage{amssymb}
\usepackage{mathtools}
\usepackage{csquotes}
\newtheorem{theorem}{Theorem}
\title{On the Divisibility Relation $\sigma(n)\mid\sigma(n+h)$ and a Generalized Erd\H{o}s--Sierpi\'nski Conjecture}
\author{Amirali Fatehizadeh, Florian Luca}
\date{Source: arXiv:2609.04980v1 (CC BY 4.0)}
\begin{document}
\maketitle

\noindent\emph{Source and licence.} On the Divisibility Relation $\sigma(n)\mid\sigma(n+h)$ and a Generalized Erd\H{o}s--Sierpi\'nski Conjecture. Version arXiv:2609.04980v1 (CC BY 4.0), distributed by arXiv under the Creative Commons Attribution 4.0 International licence, http://creativecommons.org/licenses/by/4.0/, as recorded in the arXiv licence metadata for this exact version and shown on the abstract page https://arxiv.org/abs/2609.04980v1. Full licence notice: \texttt{licenses/arxiv-2609.04980-CC-BY-4.0.txt}.\par\smallskip

\noindent\textit{Editorial scope.} Counted unit: the article's theorem thm:conditional-family (source0.tex lines 1767--1801). The article's complete original proof of it is delivered with it (source0.tex lines 1802--1911, one \texttt{proof} environment of 110 lines). No mathematics is written, repaired or re-derived: the statement and the proof are the article's own lines, in the article's own notation, and no retained line is substituted or re-flowed.  No further source-local context is needed: every cross-reference of every retained line resolves inside the two delivered spans.  Nothing was omitted from any retained span, so the package carries no omission record.  No retained line contains a citation, so the wrapper carries no bibliography and no bibliography entry.  No retained line contains a figure environment, an image-inclusion command or drawing code, so no image is delivered and the package declares no asset dependency.  Editorial formatting only, in addition: an AMS \texttt{amsart} preamble that restates the article's own declarations and macros used by the retained lines, namely the environments \texttt{theorem} on separate counters, together with the standard packages those lines need.  The retained environments print in this document as Theorem 1 in order of appearance, whereas the article prints the same items with the numbering of its own full document.  The complete native source of the article as distributed in the arXiv e-print for this exact version is bundled unchanged as \texttt{sources/209402/source0.tex}.
\begin{theorem}\label{thm:conditional-family}
The following assertions hold.

\textup{(i)} Assume Schinzel's Hypothesis $H$. Then $\mathcal T$ is infinite, and for every $t\in\mathcal T$,
\[
\sigma(N_t)=2\sigma(n_t).
\]
Moreover, the integers $n_t$, $t\in\mathcal T$, are pairwise distinct. Consequently, under Schinzel's Hypothesis $H$, the equation
\[
\sigma(n+1)=2\sigma(n)
\]
has infinitely many positive integer solutions.

\textup{(ii)} Assume the Bateman--Horn conjecture for the four linear forms $f_1,f_2,f_3,f_4$. For each prime $p$, let $\nu(p)$ denote the number of residue classes modulo $p$ on which
\[
f_1(t)f_2(t)f_3(t)f_4(t)
\]
vanishes, and put
\[
\mathfrak S
=
\prod_p\left(1-\frac1p\right)^{-4}\left(1-\frac{\nu(p)}p\right).
\]
Then $\mathfrak S>0$, and
\[
F(x)
\sim
\frac{8\mathfrak S}{3\sqrt{42}}
\frac{\sqrt{x}}{(\log x)^4}.
\]
In particular,
\[
F(x)\gg\frac{\sqrt{x}}{(\log x)^4}.
\]
\end{theorem}

\begin{proof}
We first prove part~(i). Let $t\in\mathcal T$. Since
$f_1(t)-f_2(t)=246t+218>0$, the primes $f_1(t)$ and $f_2(t)$ are
distinct. Similarly, $f_4(t)-f_3(t)=29t+25>0$, so $f_3(t)$ and
$f_4(t)$ are distinct. Since $t\in\mathcal T$, the latter two
numbers are primes exceeding $3$, and hence are coprime to $6$.
Thus multiplicativity of $\sigma$ applies to
\[
n_t=f_1(t)f_2(t),
\qquad
N_t=6f_3(t)f_4(t).
\] Thus
\[
\begin{aligned}
\sigma(n_t)
&=(252t+224)(6t+6)\\
&=28(9t+8)\cdot6(t+1)\\
&=168(t+1)(9t+8),
\end{aligned}
\]
whereas
\[
\begin{aligned}
\sigma(N_t)
&=\sigma(6)(7t+7)(36t+32)\\
&=12\cdot7(t+1)\cdot4(9t+8)\\
&=336(t+1)(9t+8).
\end{aligned}
\]
Consequently,
\begin{equation}\label{eq:multiplier-two}
\sigma(N_t)=2\sigma(n_t).
\end{equation}

We next verify the hypotheses of Schinzel's Hypothesis $H$. Each
$f_i$ is a nonconstant irreducible linear polynomial with positive
leading coefficient, and each is primitive; indeed,
$(252,223)=(6,5)=(7,6)=(36,31)=1$. For every prime $p\ge5$, the
product $f_1f_2f_3f_4$ has at most four zeros modulo $p$, and hence
fewer than $p$. For $p=2$ and $p=3$, a direct calculation shows that
the product has exactly one zero modulo $p$. Thus $f_1f_2f_3f_4$ has no
fixed prime divisor. Schinzel's Hypothesis $H$ therefore implies that $\mathcal T$ is infinite.

Finally,
\[
n_{t+1}-n_t=3024t+4110>0\qquad(t\ge0),
\]
so $t\mapsto n_t$ is strictly increasing. Hence the integers $n_t$ with $t\in\mathcal T$ are pairwise distinct. This proves part~(i).

For part~(ii), write the four forms as $a_it+b_i$. If a prime $p$ divides neither a leading coefficient $a_i$ nor any of the pairwise determinants $a_ib_j-a_jb_i$, then all four forms have well-defined roots modulo $p$, and these roots are pairwise distinct. For the present forms, the six pairwise determinants are
\[
-78,\quad-49,\quad-216,\quad1,\quad6,\quad1.
\]
Together with the prime divisors of the leading coefficients, these
show that only $2,3,7,13$ can differ from the generic four-root case. Direct calculation gives
$\nu(2)=1$, $\nu(3)=1$, $\nu(7)=2$, and $\nu(13)=3$, whereas
$\nu(p)=4$ for $p\notin\{2,3,7,13\}$.
Every local factor in $\mathfrak S$ is therefore positive. Moreover, for $p\notin\{2,3,7,13\}$,
\[
\left(1-\frac4p\right)\left(1-\frac1p\right)^{-4}
=
1-\frac6{p^2}+O(p^{-3}).
\]
Since $\sum_p\frac1{p^2}<\infty$, the product over the nonexceptional primes converges absolutely to a positive limit. The finitely many exceptional local factors are also positive. Hence $\mathfrak S>0$.

Let
\[
\Pi(T)=\#\{1\le t\le T:t\in\mathcal T\}.
\]
The Bateman--Horn conjecture gives
\[
\Pi(T)
\sim
\mathfrak S\int_2^T\frac{du}{(\log u)^4}
\sim
\mathfrak S\frac{T}{(\log T)^4}.
\]
Since
\[
n_t=1512t^2+2598t+1115,
\]
the largest integer $t$ for which $n_t\le x$ is, for all sufficiently large $x$,
\[
T_x
=
\left\lfloor
\frac{\sqrt{6048x+6084}-2598}{3024}
\right\rfloor.
\]
Thus
\[
T_x\sim\sqrt{\frac{x}{1512}},
\qquad
\log T_x=\frac12\log x+O(1).
\]
By the strict monotonicity proved above,
\[
F(x)=\Pi(T_x).
\]
It follows that
\[
\begin{aligned}
F(x)
&\sim\mathfrak S\frac{T_x}{(\log T_x)^4}\\
&\sim\frac{16\mathfrak S}{\sqrt{1512}}\frac{\sqrt{x}}{(\log x)^4}\\
&=\frac{8\mathfrak S}{3\sqrt{42}}\frac{\sqrt{x}}{(\log x)^4}.
\end{aligned}
\]
This proves part~(ii).
\end{proof}

\end{document}
